% Figure 1. Schematic invariant projection and documented mass-space connection.
\begin{figure*}[t]
\centering
\begin{tikzpicture}[x=1cm,y=1cm,line cap=round,line join=round,
  >=Stealth,every node/.style={font=\footnotesize}]
  \begin{scope}
    \node[font=\small,anchor=north west] at (-0.15,4.55) {(a)};
    \draw[line width=0.7pt,->] (0,0) -- (7.4,0)
      node[below left,font=\small] {$T_{\rm si}$};
    \draw[line width=0.7pt,->] (0,0) -- (0,4.35)
      node[above left,font=\small] {$L_{\rm si}$};
    \fill[black!12]
      (0.55,2.85) .. controls (2.1,3.55) and (3.6,3.70) .. (5.15,3.15)
      .. controls (5.55,2.95) and (5.55,2.55) .. (5.15,2.35)
      .. controls (3.6,2.90) and (2.1,2.75) .. (0.55,2.05) -- cycle;
    \fill[black!22]
      (0.70,1.15) .. controls (2.3,0.55) and (4.0,0.70) .. (5.70,1.55)
      .. controls (6.15,1.75) and (6.15,2.15) .. (5.70,2.35)
      .. controls (4.0,1.50) and (2.3,1.35) .. (0.70,1.95) -- cycle;
    \draw[line width=0.9pt,black!70]
      (0.55,2.45) .. controls (2.1,3.15) and (3.6,3.30) .. (5.15,2.75);
    \draw[line width=0.9pt,black]
      (0.70,1.55) .. controls (2.3,0.95) and (4.0,1.10) .. (5.70,1.95);
    \fill[black] (3.55,2.18) circle (2.4pt);
    \draw[line width=0.8pt,black,fill=white]
      (3.78,2.02) rectangle ++(0.28,0.28);
    \node[anchor=south] at (3.55,2.42) {A};
    \node[anchor=west] at (4.12,2.02) {B};
    \node[align=center,font=\scriptsize,anchor=north] at (3.65,0.55)
      {near-coincident images};
    \node[font=\scriptsize,anchor=west] at (0.45,3.55) {projected branch};
    \node[font=\scriptsize,anchor=west] at (0.60,0.85) {projected branch};
    \node[font=\scriptsize,anchor=north east] at (7.05,4.05) {schematic};
  \end{scope}

  \begin{scope}[shift={(8.5,0)}]
    \node[font=\small,anchor=north west] at (-0.15,4.55) {(b)};
    \draw[line width=0.7pt] (0,0) rectangle (6.6,4.0);
    \foreach \m/\lab in {0.388/0.8,2.329/0.9,4.271/1.0,6.212/1.1}
      \draw[line width=0.6pt] (\m,-0.08) -- (\m,0)
        node[below=2pt] {$\lab$};
    \foreach \m/\lab in {0.148/0.7,0.889/0.8,1.630/0.9,2.370/1.0,3.111/1.1,3.852/1.2}
      \draw[line width=0.6pt] (-0.08,\m) -- (0,\m)
        node[left=2pt] {$\lab$};
    \node[below=10pt,font=\small] at (3.3,0) {$m_1$};
    \node[rotate=90,anchor=south,font=\small] at (-0.85,2.0) {$m_2$};
    \coordinate (A) at (5.668,2.859);
    \coordinate (B) at (6.212,3.267);
    \draw[line width=1.3pt,black] (A) -- (B);
    \fill[black] (A) circle (2.6pt);
    \draw[line width=0.9pt,black,fill=white]
      (B) ++(-0.12,-0.12) rectangle ++(0.24,0.24);
    \node[anchor=east] at ($(A)+(-0.18,0.18)$) {A};
    \node[anchor=south west] at ($(B)+(0.06,0.10)$) {B};
    \node[font=\scriptsize,align=left,anchor=north west] at (0.25,3.75)
      {verified connection};
  \end{scope}
\end{tikzpicture}
\caption{\label{fig:projection}Continuation connectivity and invariant projection.
(a)~Schematic of the two branches visible in the corrected scale-invariant
period--angular-momentum representation. The documented long-range pair A and
B has nearly coincident invariant images. (b)~The same pair in the
$(m_1,m_2)$ plane. The line denotes the verified bidirectional continuation
connection between the endpoints; its drawn shape is schematic. Their mass
distance is $0.062$. The projected branch separation therefore occurs within
one continuation-connected sampled solution set.}
\end{figure*}
